% Measurement setup (owner: C). Numbers from harness/hypotheses.md and linking/protocol.md.

\section{Measurement Setup}
\label{sec:setup}

\paragraph{Logging server.} Our partner sites run on one logging HTTPS server we wrote for this study.
For every connection it records the raw TLS ClientHello (from which we derive JA3 and JA4), the ALPN
offer, and for every request the HTTP version, the ordered header list, the HTTP/2 SETTINGS, window
update and pseudo-header order (an Akamai-style fingerprint), timestamps, the client address, and any
text submitted on the page. Each visit goes to a path that carries its run identifier
(\texttt{/r/<run\_id>/}), which joins the log to a run manifest holding the ground truth: the agent, the
model and the (fictitious) user.

\paragraph{Pages and task.} Each site serves a short page with a visible request: summarize this page in
two sentences and submit it in a form, with a visible GET link as a fallback for tools that cannot
POST. Every agent receives the same fixed task in a fresh session (visit the page, summarize it,
submit). There is no hidden text and no injected instruction; the visible request stands in for what a
malicious site could obtain through prompt injection.

\paragraph{Local and shared deployments.} Tracks 1 and 2 ran on each member's own machine against a
local instance of the server (Linux, and two Windows 11 machines). A local container hides the client
address, so Track 3 runs against a shared cloud VM with two public hostnames and publicly trusted
certificates, listening on IPv4 and IPv6. On the VM only, each page view also carries one ordinary
sentence on the page's topic drawn from a pool of 20 per site, and the server logs which one it served
(Section~\ref{sec:discussion}).

\paragraph{Agents and baselines.} We study Claude Code (headless, with WebFetch only, or with WebFetch
and a shell) and the Codex CLI, on Linux and Windows. As baselines we run 12 plain HTTP stacks that the
agents' runtimes use (Node: axios in two versions, got, \texttt{https}, \texttt{fetch}; Python:
requests, httpx, aiohttp, urllib; curl; Go: net/http and Colly), each 30 times per site on each machine
and making the same requests as the task: 722 runs on Linux and 660 and 720 runs on the two Windows machines.
